\section{Discussion and limitations}
\label{sec:discussion}
\leavevmode\refstepcounter{subsection}\textbf{\thesubsection\ How many candidates, and when.} Our results qualify the conjecture of \citet{gladstone2026xm} in two ways. First, with non-repeating, identifiable contexts, a budget of $16$--$32$, not a small one such as $4$, was needed to keep memorization near chance, for 2 to 16 modes alike (at 8k steps; coverage at $m{=}16$: $0.92$--$0.99$ for 2 and 4 modes, $0.64$ and $0.28$ for 8 and 16, ideal $0.88$ and $0.64$). Second, no $m$ we tried matched the spread of the continuous, high-dimensional futures of BAIR. On the bounce geometry the needed budget did not grow with the number of modes, but we found no distortion measure that sets it quantitatively. The energy-score generator \citep{makabe2026esp} was competitive with fresh contexts or without an identifying feature and memorized otherwise.

\leavevmode\refstepcounter{subsection}\textbf{\thesubsection\ Limitations.} (1) \imlewm{} under-disperses on real video, and no variant fixed it; the isotropic fit on BAIR uses one free parameter. (2) On stochastic Push-T the advantage over flow matching is significant in precision error only. (3) The benchmarks are synthetic or simulated except for BAIR, the models are small, and the flow-matching baselines are far from the scale of \citet{porcher2026flowwm}; flow matching trained $3\times$ as long still covered fewer modes, but we did not scale the models, and we did not test consistency models \citep{song2023consistency}. (4) Planning success did not improve on either task, and we did not test a cost that counts a fall in any imagined frame. (5) Large budgets flatten mode probabilities, and the energy-score generator was not better calibrated. (6) The distortion account is qualitative: the prespecified geometry test was inconclusive, and the distortion in the loss's own metric fits the budget grid worse than a function of $m$ alone. (7) Ablations, BAIR variants and the deterministic Push-T control use one seed, and grids and most follow-up runs two. (8) We derived Proposition~\ref{prop:es} after observing the energy-score generator memorize, recorded predictions in advance only for its follow-up runs, and examined one $\beta$ and two distance choices. (9) We observed outcome memorization on one benchmark family of fixed size (26 to 38 visits per training context) and a toy problem, not on BAIR. (10) Per-step matching lacks a control at equal training time, and we did not compare with learned codebooks (Appendix~\ref{app:proofs}).

\section{Conclusion}
\label{sec:conclusion}
In our small-scale benchmarks, a one-step world model trained with the \nom{} objective ($m{=}16$) on frozen DINO tokens covered the four modes of a bounce benchmark better than both flow matching trained $3\times$ as long and the one-step flow baselines, and its rollouts stayed as close to real token maps as those of regression. At 8k steps, budgets of $16$--$32$ kept outcome memorization near chance for 2 to 16 modes, whereas the energy-score generator memorized at both budgets we tried. At fixed $m$ and $K$, memorization changed with the mode geometry, as a distortion would, but our prespecified test of its form was inconclusive. No budget up to 64 matched the spread of BAIR's futures.
